% =========================================================================
% Chapter 4: Data Description
% Sake Saakstra, MSc EOR thesis, Vrije Universiteit Amsterdam
% Supervisor: prof. dr. S.J. Koopman
% =========================================================================
\documentclass[11pt,a4paper]{article}

\usepackage[a4paper,margin=2.5cm]{geometry}
\usepackage{amsmath,amssymb,amsthm,bm}
\usepackage{graphicx}
\usepackage{booktabs}
\usepackage{natbib}
\usepackage{hyperref}
\usepackage{xcolor}
\usepackage{caption}
\usepackage{makecell}
\hypersetup{colorlinks=true,linkcolor=blue!50!black,citecolor=blue!50!black,urlcolor=blue!50!black}

\title{\textbf{Chapter 4}\\[0.3em] Data Description:\\
Constructing a Hydrogen Project Survival Panel}
\author{Sake Saakstra \\ \small Supervisor: prof. dr. S.J. Koopman}
\date{Draft v1, \today}

\newcommand{\E}{\mathbb{E}}

\begin{document}
\maketitle

\begin{abstract}
This chapter documents the data construction underpinning the empirical analyses of Chapters 5 through 7. We assemble a project-level survival panel covering 714 announced hydrogen production projects globally between 2001 and 2026, comprising 244 Blue (SMR/ATR with carbon capture) and 470 Green (PEM electrolyser) technology projects. Project-level covariates include announced capacity, region, sponsor identity and type, and cancellation status. We complement the project panel with monthly EUA spot prices over 2010--2024 (mean \euro 30, range \euro 3--96), aggregated to annual averages for matching with the project event-time structure. The resulting person-year panel contains 4{,}162 observations across 714 projects, with 43 cancellation events. We document four important data limitations: systematic under-identification of sponsors for pilot-stage projects (369 of 714 projects have ``Unknown'' sponsor), right-censoring of projects announced post-2024, geographic concentration of Blue projects in oil-major-led regions, and the absence of intermediate project-status milestones (announced $\to$ FEED $\to$ FID $\to$ construction $\to$ operational). These limitations are non-trivial for inference and are explicitly addressed in the empirical chapters.
\end{abstract}

% =========================================================================
\section{Overview and Data Philosophy}
\label{sec:ch4-overview}

The empirical analyses of this thesis require a panel of hydrogen production projects with three core ingredients: (i) a technology identifier distinguishing Blue (SMR/ATR with carbon capture and storage) from Green (PEM electrolyser-based) hydrogen production, (ii) an event indicator for project cancellation, and (iii) a contemporaneous carbon-price signal at the relevant geographical scope. The challenge of constructing such a panel is twofold. First, hydrogen project data is fragmented across industry databases, government registries, and corporate disclosures, and is not consistently maintained or audited. Second, hydrogen markets are immature: most projects are announced in the planning stage and only a minority reach financial investment decision (FID) or operational status. The distinction between ``cancelled'' and ``stalled but not formally cancelled'' is in practice blurry, and our event variable captures only formally announced cancellations.

This chapter describes our sample construction choices in detail, presents summary statistics, and documents data limitations that constrain the inferential reach of subsequent chapters.

% =========================================================================
\section{Primary Data Source}
\label{sec:ch4-source}

The primary data source is the International Energy Agency (IEA) Hydrogen Production Projects Database \citep{iea2024hydrogenprojects}, a publicly maintained registry of announced low-carbon hydrogen production facilities. The database catalogues projects from initial public announcement through to operational status, including: project name, location (country/region), announced production capacity (in MW for electrolysers or kt H$_2$/year for fossil-based plants), sponsor identity (where publicly disclosed), production technology, announced commissioning date, and current operational status as of the database's most recent update. We use the version of the database current as of January 2026.

We supplement the IEA project data with three additional sources. EU Emissions Allowance (EUA) spot price data is sourced from Refinitiv Eikon for the period January 2010 through December 2024, sampled at monthly frequency. Natural gas spot prices (TTF) and electricity prices are sourced from the same provider, used in robustness checks (Chapter 8). Project sponsor classifications --- distinguishing oil majors, industrial gas companies, utilities, pure-play hydrogen firms, steel firms, and ``other'' (including unknown) --- are constructed from the IEA sponsor field cross-referenced with corporate-registry data. We deliberately retain the granular sponsor identity (e.g., ``BP p.l.c.'', ``Equinor ASA'') alongside the type classification to enable both pooled-type and project-specific analysis.

% =========================================================================
\section{Sample Construction}
\label{sec:ch4-construction}

\subsection{Technology filter}

Our analysis focuses on the two production technologies that dominate the announced project pipeline: Blue (fossil-based reforming with carbon capture) and Green (PEM electrolyser). Together these account for over 90\% of low-carbon hydrogen project capacity announced through 2024 \citep{iea2024hydrogenprojects}.

We exclude:
\begin{itemize}
    \item Grey hydrogen projects (no carbon capture).
    \item Pure research demonstrators at sub-pilot scale ($<$ 0.1 MW).
    \item Methane pyrolysis projects (turquoise hydrogen), an emerging third pathway with insufficient sample for separate analysis.
    \item Alkaline electrolyser projects that are not specifically PEM-based, to maintain technology homogeneity within the Green category.
\end{itemize}
The exclusion of alkaline electrolyser projects is a defensible choice because PEM and alkaline have different cost structures and load-following characteristics; pooling would conflate two distinct technologies. We acknowledge in Chapter 8 that the choice limits the generality of our findings to PEM specifically rather than ``Green hydrogen broadly.''

\subsection{Time window}

The sample window spans announced projects with declared start (i.e., announcement) dates between 2001 and 2026. While the modal start year falls in 2020--2024, we retain earlier projects (which are sparse: only 33 projects announced before 2010) to capture the early-stage market context. The carbon-price signal data is available beginning 2010, so projects announced before 2010 enter the analysis with EUA assignments taking the 2010 average for pre-2010 calendar years. This conservative imputation is verified not to drive results in Chapter 8.

\subsection{Right-censoring}

A project is right-censored if, as of January 2026, it has neither been formally cancelled nor reached operational status. The largest right-censoring effect concerns projects announced in 2024--2025 with declared development durations of 4--6 years: these projects appear in the panel as ongoing observations through the end of the sample window. We address right-censoring through a partial-likelihood specification in Chapter 5 (Cox proportional hazards) and through person-year hazard models that conditionally exclude censored projects in their final observation period (Chapters 6 and 7).

% =========================================================================
\section{Variable Definitions}
\label{sec:ch4-variables}

The final v7 sample contains 714 unique projects with the following variables:

\begin{table}[h!]
\centering
\caption{Variable definitions in the v7 project panel.}
\begin{tabular}{lp{10.5cm}}
\toprule
Variable & Definition \\
\midrule
\texttt{project\_id} & Unique project identifier (0--713) \\
\texttt{is\_blue\_ccs} & 1 if Blue hydrogen (SMR/ATR with CCS); 0 if Green (PEM) \\
\texttt{tech} & String technology descriptor (Blue\_CCS or PEM) \\
\texttt{log\_capacity\_mw} & Natural log of announced production capacity in MW \\
\texttt{region} & One of: EU, Other Europe, North America, Asia, MENA, ANZ, Other \\
\texttt{sponsor\_type} & Classification: Oil\_major, Industrial\_gas, Utility, Pure\_play, Steel, Other, Unknown \\
\texttt{sponsor\_owner} & Named sponsor entity (e.g., BP plc., Equinor ASA), or ``Unknown'' \\
\texttt{year\_announced} & Calendar year of project announcement (2001--2026) \\
\texttt{duration} & Years between announcement and event (or censoring) \\
\texttt{event\_type} & 0 = no event; 1 = formal cancellation; 2 = informal abandonment \\
\texttt{event\_any} & 1 if \texttt{event\_type} $\in \{1, 2\}$; 0 otherwise \\
\bottomrule
\end{tabular}
\label{tab:ch4-variables}
\end{table}

For analytical purposes we combine \texttt{event\_type} values 1 and 2 into the binary outcome \texttt{event\_any}, treating formal cancellations and informal abandonments as equivalent failures. This pooling is conservative and is robustness-tested in Chapter 8.

The regional categorisation merges UK, Norway, and Switzerland into ``Other Europe'' because these jurisdictions have separate carbon-pricing regimes (UK ETS post-Brexit, Norwegian EUA-linked system, Swiss ETS) that are not identical to the EU ETS. The ``EU'' category captures continental EU member states under the EU ETS. The ``Asia'' category combines Japan, South Korea, China, Singapore, and India. ``ANZ'' covers Australia and New Zealand. ``MENA'' includes Middle East and North Africa, with the largest representation from Saudi Arabia and UAE. ``Other'' captures Latin America and Africa.

The carbon-price signal $z_t$ used in the empirical models is constructed as follows. We compute monthly average EUA spot price from the daily series and aggregate to annual averages. The annual EUA series is then standardised by subtracting the sample mean and dividing by the sample standard deviation:
\begin{equation}
z_t = \frac{\mathrm{EUA}_t - \bar{\mathrm{EUA}}}{\mathrm{SD}(\mathrm{EUA})},
\end{equation}
where $\bar{\mathrm{EUA}} = $ \euro 30 and $\mathrm{SD}(\mathrm{EUA}) = $ \euro 29 over the 2010--2024 period. The standardised $z_t$ has mean zero and unit variance in the empirical sample, with extreme values reaching $z = +2.3$ in the 2023 EUA peak and $z = -0.9$ in the 2010--2017 trough.

% =========================================================================
\section{Summary Statistics}
\label{sec:ch4-summary}

\subsection{Project distribution by technology and region}

Table \ref{tab:ch4-region-tech} cross-tabulates project counts by region and technology. The 714 projects in our sample are concentrated in EU continental (213 projects, 30\% of sample), North America (162 projects, 23\%), Asia (106 projects, 15\%), and Other Europe (87 projects, 12\%). PEM projects are particularly concentrated in the EU (193 of 470 PEM projects, 41\%), reflecting both the EU's Hydrogen Strategy (2020) and the Innovation Fund's targeting of electrolyser deployment.

\begin{table}[h!]
\centering
\caption{Project counts by region and technology in v7 sample.}
\begin{tabular}{lccc}
\toprule
Region & PEM & Blue & Total \\
\midrule
EU                  & 193 &  20 & 213 \\
Other Europe        &  54 &  33 &  87 \\
North America       &  64 &  98 & 162 \\
Asia                &  66 &  40 & 106 \\
MENA                &   7 &  12 &  19 \\
ANZ                 &  37 &  14 &  51 \\
Other               &  49 &  27 &  76 \\
\midrule
Total               & 470 & 244 & 714 \\
\bottomrule
\end{tabular}
\label{tab:ch4-region-tech}
\end{table}

Blue projects are conversely concentrated in North America (98 of 244 Blue projects, 40\%), where oil-major sponsorship dominates. This geographic distribution of technology is itself substantively informative: it reflects the comparative advantage of regions with mature gas infrastructure (North America, MENA) for Blue, versus regions with strong renewable electricity capacity (EU, ANZ) for Green.

\subsection{Capacity and duration}

Median announced capacity is 1.0 MW for Blue projects and 5.0 MW for PEM projects, suggesting that the typical Blue facility in the announcement pipeline is smaller in scale-equivalent than the typical PEM facility. Mean (log) capacity is 0.12 for Blue and 2.13 for PEM, with PEM exhibiting a long right tail of large projects (up to 1,000 MW gigafactory-scale announcements). Mean project duration --- from announcement to event or censoring --- is 4.25 years for Blue and 5.33 years for PEM, reflecting both technological maturity (Blue's reliance on established SMR/ATR designs allows faster timelines when committed) and project commercialisation status.

\subsection{Sponsor distribution}

The sponsor classification distribution underscores the data-quality challenges that motivate much of our empirical methodology. Of 714 projects, 369 (52\%) have sponsor identity recorded as ``Unknown.'' Among the 345 projects with identified sponsors, the breakdown is: Oil majors (24 projects, e.g., BP, Equinor, Shell), Industrial gas companies (19, e.g., Linde, L'Air Liquide), Utilities (6), Pure-play hydrogen companies (3), Steel producers (2), and Other identified entities (291, including renewable developers, port authorities, methanol producers).

Strikingly, the ``Unknown'' category includes 369 PEM projects but only a handful of Blue projects, reflecting the latter's dominance by major energy companies with public disclosure. This systematic correlation between technology and sponsor-identification status emerges as a confounding factor in our subsequent empirical analysis (Chapter 7) and is explicitly controlled for in robustness specifications.

\subsection{Event distribution}

The v7 sample contains 43 cancellation events distributed as 27 Blue and 16 Green. The temporal distribution of events is highly concentrated: 20 events occurred in 2023 alone (47\% of all events), and a further 14 in 2024 (33\%). Combined, the 2023--2024 period accounts for 79\% of all observed cancellations. This concentration reflects three confluences: (i) the post-2022 energy-crisis correction of EUA prices and gas-price expectations; (ii) the maturation of projects announced during the 2020--2021 hydrogen-strategy boom; and (iii) the cumulative budgetary pressure on sponsors of multi-year capital-intensive projects.

\begin{table}[h!]
\centering
\caption{Cancellation events by year and technology.}
\begin{tabular}{lccc}
\toprule
Event Year & Blue & PEM & Total \\
\midrule
2018 & 3 & 0 & 3 \\
2020 & 1 & 1 & 2 \\
2022 & 2 & 0 & 2 \\
2023 & 9 & 11 & 20 \\
2024 & 11 & 3 & 14 \\
2025 & 1 & 0 & 1 \\
2026 & 0 & 1 & 1 \\
\midrule
Total & 27 & 16 & 43 \\
\bottomrule
\end{tabular}
\label{tab:ch4-events-year}
\end{table}

The aggregate cancellation rate over the sample period is 11\% for Blue projects (27 of 244) and 3.4\% for PEM projects (16 of 470). This raw ratio --- a Blue-to-PEM cancellation rate differential of approximately 3.2 --- is consistent with the negative carbon-conditional finding documented in Chapters 6 and 7, but the unconditional rate confounds the technology-specific carbon sensitivity with both regional and sponsor-composition effects that the formal models address.

\subsection{EUA price distribution}

The EUA spot price over 2010--2024 exhibits substantial variation, with a mean of \euro 30, standard deviation of \euro 29, and range of \euro 3 (early-2010s trough) to \euro 96 (early-2023 peak). The price is highly autocorrelated; we use annual averages to align with the project-level event timing structure. The post-2017 regime (mean \euro 39, SD \euro 26) is substantively different from the pre-2017 regime (mean \euro 9, SD \euro 4), reflecting the Market Stability Reserve activation in 2019 and the subsequent supply-tightening induced by the energy crisis.

% =========================================================================
\section{Panel Construction for Hazard Modelling}
\label{sec:ch4-panel}

\subsection{Person-year panel}

For hazard estimation we convert the project-level data into a person-year panel. For each project $i$, we generate one observation for each calendar year between the announcement year $T^{\mathrm{start}}_i$ and the event/censoring year $T^{\mathrm{end}}_i = T^{\mathrm{start}}_i + \mathrm{duration}_i$. Each row in the panel includes: project identifier, calendar year, technology indicator, sponsor identity, log capacity, year-since-announcement, EUA at the calendar year, and event indicator (1 if event occurred in that year, 0 otherwise).

The resulting panel contains 4,162 person-year observations across 714 projects. Of these, 43 person-year observations contain the event indicator equal to 1 (one per project). The annualised EUA-by-year is merged via calendar year, with the standardised carbon-price signal $z_t$ replacing the raw EUA value for analytical convenience.

\subsection{Aggregated year-by-technology panel}

For the time-varying parameter analysis of Chapter 7 we further aggregate the person-year panel to a year-by-technology level: for each combination of (calendar year, technology), we compute the count of at-risk observations and the count of events. This produces a Binomial-format panel with $2 \times 17 = 34$ observations (Blue and Green at each year from 2010 through 2026), suitable for parameter-driven and observation-driven time-varying parameter specifications.

The aggregated panel preserves the unconditional hazard information while sacrificing the individual-project covariates. This trade-off is appropriate because the time-varying parameter analysis is primarily concerned with the time-evolution of the conditional carbon coefficient $\beta_{\mathrm{int}}(t)$, for which sample-mean residual covariates are sufficient.

% =========================================================================
\section{Data Limitations}
\label{sec:ch4-limitations}

Four limitations of the v7 sample warrant explicit acknowledgement and are returned to in the discussion of Chapter 10.

\subsection{Sponsor under-identification}

The 369 projects with sponsor recorded as ``Unknown'' constitute 52\% of the sample. These are systematically smaller projects, more likely PEM than Blue (Section \ref{sec:ch4-summary}), and concentrated in pilot/demonstration stages. The under-identification reflects a combination of (i) genuine non-disclosure by sponsors, (ii) database curation gaps, (iii) consortium-based projects with no single legal owner, and (iv) academic/research-stage projects that have not been propagated through commercial registries. We address this limitation in Chapter 7 through sponsor-controls in the regression specification and through stratified analysis by sponsor-known status.

\subsection{Right-censoring of post-2024 projects}

Projects announced in 2024--2025 with declared 4--6 year development durations are largely right-censored in our sample. Of 144 projects announced in 2024--2026, only a handful have reached operational status; the remainder are mid-development. This censoring is informative if our econometric specifications are robust to it (and we verify this in Chapter 5 via Cox PH with explicit censoring); it would be problematic if there were systematic differences between projects that have already cancelled and projects still in development.

\subsection{Geographic concentration}

Blue projects are heavily concentrated in North America (40\%) and Other Europe (UK and Norway, 14\%). The European Union proper contains only 20 Blue projects of 244, complicating the test of EUA price effects in the most direct ETS-bound regime. We address this limitation in Chapter 7 through regional stratification and by acknowledging that the cleanest available signal comes from North American data.

\subsection{Outcome definition}

Our event variable captures only formally announced cancellations and informal abandonments. Several intermediate project failures are not captured:
\begin{itemize}
    \item Projects that pass FEED but never reach FID (effectively cancelled but never formally announced).
    \item Projects that downsize substantially from original announcement (partial failure).
    \item Projects that switch technology mid-development (e.g., from Blue to Green).
    \item Projects that experience prolonged delays without formal status change.
\end{itemize}
The omission biases the estimated cancellation rate downward, especially for projects with sponsor identities that do not publicly announce setbacks. This is more pronounced for Unknown-sponsor pilots than for major-sponsor industrial projects.

% =========================================================================
\section{Conclusion}
\label{sec:ch4-conclusion}

The v7 sample assembled in this chapter forms the empirical foundation of Chapters 5 through 7. The sample comprises 714 announced hydrogen production projects (244 Blue, 470 Green) globally over 2001--2026, with 43 formal cancellation events and a person-year panel of 4{,}162 observations. The data is supplemented by EUA spot prices and sponsor classifications. Four data limitations --- sponsor under-identification, right-censoring of recent projects, geographic concentration of Blue projects in North America, and the formal-announcement definition of cancellation --- are non-trivial and are explicitly accounted for in subsequent empirical chapters. Despite these limitations, the v7 sample is, to the best of our knowledge, the largest and most comprehensive panel of hydrogen project survival data assembled for academic econometric analysis to date.

% =========================================================================
\bibliographystyle{apalike}
\begin{thebibliography}{99}

\bibitem[IEA(2024)]{iea2024hydrogenprojects}
International Energy Agency (2024).
\newblock \textit{Hydrogen Production Projects Database, 2024 Update}.
\newblock IEA, Paris.

\bibitem[IEA(2024b)]{iea2024hydrogen}
International Energy Agency (2024).
\newblock \textit{Global Hydrogen Review 2024}.
\newblock IEA, Paris.

\bibitem[Odenweller et~al.(2022)]{odenweller2022probabilistic}
Odenweller, A., Ueckerdt, F., Nemet, G.~F., Jensterle, M., \& Luderer, G. (2022).
\newblock Probabilistic feasibility space of scaling up green hydrogen supply.
\newblock \textit{Nature Energy}, 7(9), 854--865.

\end{thebibliography}

\end{document}
